\scanpage{42}
% slanted
\margin{\uline{à vérifier}}
\uline{est contenue dans une réunion d'images} \uline{de parties} $M_K$-\uline{bornées de} $U_i(\overline K)\times M_{\overline K}$.

Admettons ce point. On aura, par l'\uline{exemple précédent}, un hom de faisceaux
\diagram{\begin{tikzcd}[ampersand replacement=\&]
\underline{O}_X^*\arrow[r,"{\substack{f\mapsto\varphi_f\\ \deleted{\unsure{\xi|f}}}}"] \& \underline{\Phi}_X
\end{tikzcd}}
d'autre part on a un hom de faisceaux
\diagram{\begin{tikzcd}[ampersand replacement=\&]
{\deleted{R}\unsure{O|\underline O|\underline{\mathcal O}}}_X^*\arrow[r,hook,"{\substack{\unsure{f\leftarrow\!\vert f|f\mapsto f}\\ \deleted{\xi}}}"] \& \underline{R}_X^{\times}
\end{tikzcd}}
où $R_X^*$ est le faisceau des fonctions rationnelles sur $X$ [N.B. nous supposons $X$ réduit ou $=$ séparable $/K$, pour fixer les idées]. Nous voulons ``une généralisation commune des fonctions admissibles et des fonctions rationnelles'', et définissons pour ceci le \uline{faisceau des quasi-fonctions} comme le faisceau
\[
\begin{aligned}
\underline{\mathrm{QF}}_X&=\underline{\Phi}_X\mathop{\tikz[baseline=(op.base)]{\node[inner sep=0pt] (op) {$\unsure{\amalg|\times}$};\draw[overlay,thin] (-.33,-.22) .. controls (.09,-.33) and (-.04,.48) .. (.35,.49);}}\limits_{O_X^*}\underline R_X^{\times}
&&\text{(somme amalgamée)}\\
&\simeq\underline{\Phi}_X\mathbin{\tikz[baseline=(op.base)]{\node[inner sep=0pt] (op) {$\times$};\draw[overlay,thin] (-.33,-.22) .. controls (.09,-.33) and (-.04,.48) .. (.35,.49);}}\underline R_X^*/\unsure{\Theta|\theta}(\underline O_X^*)
\end{aligned}
\]
% archive sheet number: 41
